\chapter{DevSecOps}
\label{chap:devsecops}

\section{Containerizzazione}
\label{sec:containerizzazione}

Tutti i microservizi del sistema sono containerizzati con \textbf{Docker}, usando
immagini base minimali per ridurre la superficie di attacco e le dimensioni dell'immagine:

\begin{table}[ht]
\centering
\small
\begin{tabularx}{\linewidth}{l l l}
\toprule
\textbf{Servizio} & \textbf{Immagine base} & \textbf{Note} \\
\midrule
backend       & \texttt{python:3.9-slim}  & Uvicorn + FastAPI \\
email\_poller & \texttt{python:3.14-slim} & APScheduler + IMAPClient \\
agent\_worker & \texttt{python:3.12-slim} & Atomic Agents + sentence-transformers \\
mcp\_server   & \texttt{python:3.12-slim} & FastMCP + SQLAlchemy \\
frontend      & Multi-stage: \texttt{node:alpine} + \texttt{nginx:alpine} & Build + serve statico \\
postgres      & \texttt{postgres:16}      & Immagine ufficiale \\
tika          & \texttt{apache/tika:latest} & Servizio REST \\
\bottomrule
\end{tabularx}
\caption{Immagini base usate per la containerizzazione.}
\label{tab:docker_images}
\end{table}

Il Dockerfile del frontend usa una \textbf{build multi-stage}:

\begin{lstlisting}[language=bash, caption={Struttura del Dockerfile multi-stage del frontend.}]
# Stage 1: build con Node
FROM node:alpine AS builder
WORKDIR /app
COPY package*.json ./
RUN npm ci
COPY . .
ARG VITE_API_URL
ENV VITE_API_URL=$VITE_API_URL
RUN npm run build

# Stage 2: serve con Nginx
FROM nginx:alpine
COPY --from=builder /app/dist /usr/share/nginx/html
EXPOSE 80
\end{lstlisting}

Questo pattern riduce l'immagine finale a soli file statici serviti da Nginx,
senza includere Node.js, npm e le dipendenze di sviluppo.

\section{Pipeline CI/CD}
\label{sec:cicd}

La pipeline CI/CD \`e implementata con \textbf{GitHub Actions} nel file
\texttt{.github/workflows/docker-build-push.yml}.

\paragraph{Trigger.} Push su branch \texttt{main}.

\paragraph{Job paralleli.} Cinque job indipendenti costruiscono e pubblicano
le immagini in parallelo: \texttt{backend}, \texttt{frontend}, \texttt{email\_poller},
\texttt{agent\_worker}, \texttt{mcp\_server}.

\paragraph{Registry.} \textbf{GitHub Container Registry} (\texttt{ghcr.io}):
\begin{itemize}[leftmargin=*]
    \item \texttt{ghcr.io/matteociaroni/dms-backend:latest}
    \item \texttt{ghcr.io/matteociaroni/dms-frontend:latest}
    \item \texttt{ghcr.io/matteociaroni/dms-email\_poller:latest}
    \item \texttt{ghcr.io/matteociaroni/dms-agent\_worker:latest}
    \item \texttt{ghcr.io/matteociaroni/dms-mcp\_server:latest}
\end{itemize}

\paragraph{Multi-platform.} Ogni immagine \`e compilata per \texttt{linux/amd64} e
\texttt{linux/arm64} tramite QEMU e Docker Buildx, garantendo compatibilit\`a con
macchine Apple Silicon e server ARM cloud.

\paragraph{Layer caching.} Il backend di cache \texttt{type=gha} (GitHub Actions)
riduce il tempo medio di build a meno di 3~minuti per immagine, amortizzando il
costo di reinstallazione delle dipendenze Python.

\begin{lstlisting}[language=yaml, caption={Job CI/CD per il backend (\texttt{.github/workflows/docker-build-push.yml}).}]
backend:
  runs-on: ubuntu-latest
  permissions:
    contents: read
    packages: write
  steps:
    - uses: actions/checkout@v4
    - uses: docker/login-action@v3
      with:
        registry: ghcr.io
        username: ${{ github.actor }}
        password: ${{ secrets.GITHUB_TOKEN }}
    - uses: docker/setup-qemu-action@v3
    - uses: docker/setup-buildx-action@v3
    - uses: docker/build-push-action@v6
      with:
        context: ./backend/
        push: true
        tags: ghcr.io/matteociaroni/dms-backend:latest
        cache-from: type=gha
        cache-to: type=gha,mode=max
        platforms: linux/amd64,linux/arm64
\end{lstlisting}

\section{Gestione dei segreti}
\label{sec:segreti}

Tutti i segreti sono iniettati tramite \textbf{variabili d'ambiente}, mai hardcoded
nel codice sorgente:

\begin{itemize}[leftmargin=*]
    \item \texttt{EMAIL\_ENCRYPTION\_KEY}: chiave Fernet per la cifratura delle
          credenziali IMAP. Generata con:
\begin{lstlisting}[language=bash]
python -c "from cryptography.fernet import Fernet; print(Fernet.generate_key().decode())"
\end{lstlisting}
    \item \texttt{CUSTOM\_API\_KEY}: API key del provider LLM, passata solo
          all'agent worker.
    \item \texttt{POSTGRES\_PASSWORD}, \texttt{S3\_ACCESS\_KEY},
          \texttt{S3\_SECRET\_KEY}: credenziali infrastrutturali.
\end{itemize}

In produzione, le variabili d'ambiente devono essere gestite tramite un secrets manager
dedicato (Kubernetes Secrets con cifratura a riposo, HashiCorp Vault, AWS Secrets Manager).

\section{Schema DB versionato}
\label{sec:db_versioning}

Lo schema del database \`e definito in \texttt{init.sql} ed eseguito automaticamente
da PostgreSQL al primo avvio tramite il meccanismo
\texttt{/docker-entrypoint-initdb.d/}. Il file \`e mount-ato in sola lettura nel container.

Per l'evoluzione dello schema in produzione, lo strumento raccomandato \`e
\textbf{Alembic} (migration framework per SQLAlchemy), che permetterebbe di applicare
incrementalmente le modifiche senza perdita di dati. \TODO{Integrare Alembic per la
gestione delle migration incrementali}

\section{Deployment su Kubernetes}
\label{sec:kubernetes}

\TODO{Questa sezione descriver\`a i manifest Kubernetes per il deployment in produzione,
comprendendo: Deployment e ReplicaSet per i servizi stateless (backend, email\_poller,
agent\_worker, mcp\_server, frontend), StatefulSet per PostgreSQL con
PersistentVolumeClaim, Deployment per SeaweedFS con storage class appropriato,
Service per l'esposizione interna, Ingress con TLS termination per il frontend e il backend,
HorizontalPodAutoscaler per backend e agent\_worker, ConfigMap per le variabili
non sensibili e Secret per quelle sensibili.}

\section{Aree di miglioramento}
\label{sec:miglioramenti}

La revisione del codice ha evidenziato le seguenti aree migliorabili, documentate
per trasparenza:

\begin{itemize}[leftmargin=*]
    \item \textbf{CORS permissivo}: \texttt{allow\_origins=["*"]} va ristretto al
          dominio del frontend in produzione.
    \item \textbf{Uvicorn reload in compose}: la flag \texttt{--reload} del comando
          uvicorn \`e appropriata per lo sviluppo ma deve essere rimossa in produzione.
    \item \textbf{Rate limiting}: non implementato. In produzione va aggiunto tramite
          \texttt{slowapi} o a livello di Ingress/API gateway.
    \item \textbf{Credenziali S3 di esempio}: il file \texttt{s3.json} contiene
          credenziali placeholder; il deployment reale deve generarle e ruotarle
          periodicamente.
    \item \textbf{OAuth2 per email}: attualmente supportati solo \texttt{app\_password};
          il flusso OAuth2 Google/Outlook per i token \`e modellato nello schema ma
          non completamente implementato.
\end{itemize}
